% =============================================================================
% §1.9 La convergence sur le réseau (1975–2007)
% =============================================================================
% Résultats dimensionnels visés :
%   0-TER  — la numérisation crée une consommation permanente (STORE actif)
%   D2     — plateformes et marchés winner-take-all
%   D3     — capital intangible ≥ tangible (Corrado et al. 2005)
%   D7     — coûts de transaction quasi nuls, modularité (Langlois 2003)
%   D8     — cycle Perez (dotcom = railway mania)
%   D10    — open standards comme dispositif de coordination
% Sources empiriques principales :
%   Hilbert & López 2011, David 1990, Corrado/Hulten/Sichel 2005/2009,
%   Perez 2002, Koomey 2011, Flamm 1988, Shapiro & Varian 1999
% =============================================================================
 
\section{La convergence sur le réseau (1975--2007)}\label{sec:convergence}
 
La section précédente s'achevait sur une concomitance ignorée : en 1974, le microprocesseur, le \emph{Dennard scaling} et le premier choc pétrolier coexistent sans que les trajectoires du calcul, de l'énergie et de l'économie se croisent. La période qui s'ouvre ici est celle où les trois fonctions --- transmettre, stocker, calculer --- se numérisent sur le même substrat électrique, mais pas au même rythme, pas selon les mêmes logiques, et pas avec les mêmes conséquences.
 
Ce que la mesure empirique va confirmer, la théorie des systèmes sociotechniques l'avait anticipé. \textcite{ropohl_systemtheorie_1979}, dans son~\emph{Allgemeine Technologie} dont la première édition paraît en 1979, observe~: \og~Le téléphone, le télex, la radio, la télévision et l'ordinateur ont jusqu'à présent constitué des mondes séparés, aussi bien dans leurs réseaux que dans leurs terminaux. Mais ces mondes distincts convergent désormais vers un cosmos unifié~--- un réseau de communication mondial intégré, avec des terminaux universels capables de recevoir, traiter et envoyer indifféremment du texte, du son, de l'image et du film.~\fg\footnote{\og~Bislang sind Telefon, Fernschreiben, Rundfunk, Fernsehen und Computer sowohl mit ihren Netzen wie auch mit ihren Endgeräten getrennte Welten gewesen. Inzwischen aber wachsen diese verschiedenen Welten zu einem umfassenden Kosmos zusammen, zu einem globalen integrierten Kommunikationsnetz mit universellen Endgeräten, die gleichermassen Schrift, Ton, Bild und Film empfangen, verarbeiten und versenden können.~\fg~\parencite[Fallbeispiel \og~Computer~\fg, 3.~Aufl.]{ropohl_systemtheorie_1979}. L'observation est formulée en~1979, au moment précis où TCP/IP est en cours de développement (il sera adopté en 1983) et où le Web n'existe pas encore. Ropohl la confirme dans la troisième édition de 2009, ajoutant uniquement la référence à la \og~loi de Moore~\fg{} et à l'intégration ordinateur-téléphone mobile, qu'il décrit comme \og~en cours~\fg{} à la date de révision.} L'observation est formulée à un moment où la convergence n'est pas encore réalisée~--- TCP/IP sera adopté en 1983, le Web inventé en 1989~--- mais où elle est déjà inscrite dans la logique des systèmes sociotechniques existants. Ce que Ropohl appelle \emph{informationstechnische Integration} est exactement ce que la période 1975--2007 produit empiriquement.

\dimmark{0-TER}
L'étude de référence pour saisir cette période dans sa globalité est celle de \textcite{hilbert_worlds_2011}, publiée dans \emph{Science}, qui mesure les capacités technologiques mondiales de stockage, de communication et de calcul en suivant soixante technologies analogiques et numériques entre 1986 et 2007. L'étude repose sur plus de mille cent sources primaires, documentées dans un supplément méthodologique de trois cent deux pages \parencite{lopez_methodological_2012}. Trois résultats structurent la période.
 
Le premier est l'asymétrie des taux de croissance. La capacité de calcul à usage général croît de 58~\% par an entre 1986 et 2007, soit un doublement tous les dix-huit mois et une vitesse dix fois supérieure à celle du PIB mondial. La capacité de télécommunication bidirectionnelle croît de 28~\% par an, cinq fois plus vite que le PIB. La capacité de stockage installée croît de 23~\% par an, quatre fois plus vite. La diffusion de l'information par broadcasting, en revanche, ne croît que de 6~\% par an, au rythme de l'économie. Le calcul croît donc dix fois plus vite que le stockage, qui croît lui-même quatre fois plus vite que l'économie : les trois fonctions ne sont pas simplement en expansion ; elles divergent.
 
Le deuxième résultat est l'asynchronie de la numérisation. Les télécommunications sont numériques dès 1990 (99,9~\% en 2007). Le stockage ne bascule qu'au début des années 2000 (25~\% numérique en 2000, 94~\% en 2007 ; les cassettes VHS analogiques dominaient encore en 1986). Le broadcasting reste à 75~\% analogique en 2007 (la télévision hertzienne terrestre n'est éteinte aux États-Unis qu'en 2009, en France en 2011). La numérisation n'est pas un événement unique ; c'est un processus de trente ans dont le rythme diffère par fonction.
 
Le troisième résultat est la composition des fonctions dans le substrat numérique. \textcite{hilbert_worlds_2011} définissent le calcul, à la suite de \textcite{turing_computable_1936}, comme « la transmission répétée d'information à travers l'espace (communication) et le temps (stockage), guidée par une procédure algorithmique ». Autrement dit, le calcul numérique contient structurellement de la transmission et du stockage. Lorsque les trois fonctions partagent le même substrat, leurs demandes énergétiques se composent au lieu de rester séparées. C'est dans cette composition que se loge le couplage structurel entre information et énergie que les cas précédents n'avaient fait qu'entrevoir.
 
 
% -----------------------------------------------------------------------------
\subsection{Le calcul distribué sur infrastructure analogique (1975--1995)}\label{subsec:pc}
% -----------------------------------------------------------------------------
 
\dimmark{D10}
En 1975, le microprocesseur Intel 8080 permet à Ed Roberts de construire l'Altair 8800, vendu en kit à 439 dollars par correspondance. Deux ans plus tard, l'Apple~II et le TRS-80 rendent l'ordinateur personnel accessible à un public non spécialiste. En 1981, IBM entre sur le marché avec un ordinateur personnel fondé sur des standards ouverts : le processeur Intel 8088, le système d'exploitation MS-DOS de Microsoft, et une architecture documentée que n'importe quel fabricant peut reproduire. La décision d'IBM de publier les spécifications de sa machine, prise sous la pression d'un développement accéléré (le Project Chess ne dure qu'un an), crée les conditions d'un marché concurrentiel par clonage. En trois ans, Compaq, Dell et des dizaines de fabricants asiatiques produisent des compatibles IBM-PC à des prix décroissants. \textcite{vanderkooij_information_2023} identifie cet épisode comme le déclencheur de la Third Computer Revolution (1980--2000).

\dimmark{D5/D10}{+}%
La pénétration du PC dans les foyers américains, négligeable au début de la décennie, atteint quinze pour cent en 1989~; sur la même période, le nombre de \emph{network hosts} connectés aux réseaux interuniversitaires passe d'environ deux mille en 1985 à cent soixante mille à la fin de la décennie. Cette diffusion ne tient pas qu'à la trajectoire technique de l'Altair vers l'IBM~PC~; elle est médiatisée par une culture qui se forme dans la baie de San Francisco entre 1968 et 1985 et qui infléchit durablement l'architecture des systèmes. Le \emph{Whole Earth Catalog} de Stewart Brand, publié pour la première fois en 1968, articule une critique technologique de contre-culture avec une attention aux outils «~de petite échelle~» dont les ordinateurs personnels feront partie~\parencite{turner_counterculture_2008, kirk_whole_2001}. Ted Nelson publie en 1974 \emph{Computer Lib~/ Dream Machines}, manifeste auto-édité qui appelle les utilisateurs à défier la «~\emph{computer priesthood}~» d'IBM et du gouvernement et à concevoir les ordinateurs comme \emph{creativity systems} plutôt que comme machines de calcul d'entreprise~\parencite{nelson_computer_1974}. La trajectoire des fondateurs d'Apple, Steve Jobs et Steve Wozniak, croise celle du \emph{Homebrew Computer Club} de Menlo Park où se diffusent ces valeurs~\parencite{markoff_dormouse_2005, levy_hackers_1984}.

\dimmark{D2/0-GOV}{*}%
La tension entre la culture du partage gratuit et la nécessité de monétisation de l'industrie naissante se cristallise au tournant de la décennie. Le 3~février 1976, Bill Gates~--- alors âgé de vingt ans et cofondateur d'une jeune entreprise nommée Micro-Soft~--- publie dans le bulletin MITS \emph{Altair Computer Notes} une «~lettre ouverte aux hobbyists~» dans laquelle il reproche aux utilisateurs de copier sans autorisation l'interpréteur BASIC qu'il a écrit pour l'Altair~\parencite{gates_hobbyists_1976}. L'argument est économique~: si le logiciel ne se vend pas, personne ne pourra investir dans son développement. La réponse de la communauté est largement hostile~: la pratique du partage est tenue pour intrinsèque à la pratique informatique elle-même, et l'idée de \emph{vendre} du logiciel apparaît comme une trahison des valeurs partagées. La résolution institutionnelle de cette tension se fera dans deux directions concurrentes que la suite du chapitre suivra~: le mouvement \emph{Free Software} fondé par Richard Stallman en 1984~\parencite{stallman_gnu_1984}, qui codifie la \emph{copyleft license} comme architecture juridique alternative et donnera GNU puis Linux~; et le modèle de licensing massif que Gates et Allen construisent à partir de 1980 et qui fera de Microsoft l'entreprise la plus valorisée du monde en 1998~\parencite{campbell_kelly_computer_2014}. Les deux architectures coexistent encore aujourd'hui dans un partage fonctionnel~: les systèmes dérivés d'Unix libre dominent les serveurs et les infrastructures cloud, Microsoft domine les postes de travail et le logiciel d'entreprise. \textcite{johnstone_deep_2026} identifient cette persistance comme une trace structurelle de la phase \emph{counterculture} dans l'architecture même du système informationnel contemporain, qui contraint encore la trajectoire un demi-siècle après sa cristallisation.

\dimmark{D7}
Le tableur transforme la fonction économique du micro-ordinateur. VisiCalc (1979), conçu par Dan Bricklin et Bob Frankston pour l'Apple~II, est la première application qui justifie l'achat d'un ordinateur par une entreprise pour des raisons de productivité et non de prestige technologique. Lotus 1-2-3 (1983) prend le relais sur l'IBM PC. Le tableur ne remplace pas un employé ; il modifie la manière dont les cadres interagissent avec les données financières, permettant la simulation instantanée de scénarios. Le parallèle avec la tabulatrice au Nautical Almanac Office (§\ref{subsec:hollerith}) est exact : la machine ne supprime pas le travail, elle en change la nature et la localisation.
 
\dimmark{D4}
Pendant cette phase, le calcul est distribué mais isolé. Le micro-ordinateur ne communique que par le réseau téléphonique analogique, via un modem dont le débit ne dépasse pas 14\,400 bits par seconde en 1995. Le stockage est local : disquettes, puis disques durs dont la capacité passe de 10 mégaoctets (IBM PC XT, 1983) à quelques gigaoctets au milieu des années 1990. L'infrastructure de communication reste analogique ; le réseau téléphonique commuté (PSTN) porte des données numériques sous forme de signaux modulés. \textcite{hilbert_worlds_2011} montrent que la capacité de télécommunication par habitant ne passe que de 0,16 à 0,23 mégaoctet compressé par jour entre 1986 et 1993, un doublement bien plus lent que celui de la capacité de calcul (de 0,06 à 0,8 MIPS par habitant, soit un facteur treize sur la même période).
 
\dimmark{D1}
C'est dans ce décalage entre la puissance de calcul locale et la bande passante de communication que s'installe le paradoxe de la productivité. En 1987, Robert Solow observe que « l'on voit l'ère de l'ordinateur partout sauf dans les statistiques de productivité ». \textcite{david_dynamo_1990}, dans un article court mais décisif publié dans l'\emph{American Economic Review}, propose une explication par analogie historique : l'électrification avait suivi le même schéma. L'ampoule est inventée en 1879, mais en 1900, seuls 3~\% des foyers américains disposent de l'éclairage électrique et les moteurs électriques ne représentent que 5~\% de la force motrice industrielle. Il faut attendre les années 1920 pour que la diffusion atteigne 50~\% et que les gains de productivité deviennent mesurables. David montre que le délai s'explique par le coût de remplacement d'installations encore fonctionnelles et par la nécessité de réinventer l'organisation du travail (le passage du moteur central aux moteurs individuels par machine).
 
Le parallèle dynamo/computer est le point de départ méthodologique de cette thèse (§\ref{sec:electrification}). Ce que David ne pose pas, et que les données de \textcite{hilbert_worlds_2011} permettent de voir, c'est que le décalage entre calcul et communication n'est pas seulement un problème de diffusion organisationnelle ; c'est un décalage physique entre deux taux de numérisation. Le calcul est numérique depuis 1945 ; la communication ne le devient qu'en 1990. Pendant quarante-cinq ans, le calcul numérique circule sur un réseau analogique. Le paradoxe de la productivité est, entre autres, le symptôme de cette hétérogénéité de substrat.
 
\emergence{Le PC distribue le COMPUTE chez l'utilisateur, mais le TRANSMIT reste analogique et le STORE reste local. La consommation d'énergie du traitement de l'information est encore ponctuelle : l'ordinateur peut être éteint, la disquette ne consomme rien sur l'étagère. Le profil énergétique de cette phase est \emph{additif} : le PC s'ajoute au papier sans le remplacer (la consommation de papier de bureau augmente dans les années 1980--1990).}
 
 
% -----------------------------------------------------------------------------
\subsection{La numérisation du réseau et la double infrastructure (1983--2000)}\label{subsec:www}
% -----------------------------------------------------------------------------
 
\dimmark{D10}
L'adoption de TCP/IP comme standard d'Internet en janvier 1983 est le moment technique où la transmission d'information bascule dans le numérique. La commutation par paquets transforme la nature du TRANSMIT : dans le réseau téléphonique commuté, chaque appel monopolise un circuit physique de bout en bout ; dans le réseau à paquets, l'information est découpée, routée par des algorithmes (ce qui est du calcul), mise en cache dans des mémoires tampons (ce qui est du stockage), et réassemblée à destination. Le TRANSMIT numérique n'est plus une fonction pure ; il contient structurellement du COMPUTE et du STORE. Les données de \textcite{hilbert_worlds_2011} confirment cette transformation : la capacité de télécommunication effective passe de 0,23 mégaoctet compressé par habitant et par jour en 1993 à 1,01 en 2000, puis à 27 en 2007 --- soit une multiplication par cent dix-sept en quatorze ans, tirée par l'introduction du haut débit.
 
\dimmark{0-TER}
La transition vers le numérique s'effectue cependant \emph{sur} l'infrastructure analogique existante, pas à sa place. Le modem dial-up des années 1990 transmet des paquets TCP/IP sur la boucle locale en cuivre du réseau téléphonique. Le DSL, introduit à la fin de la décennie, réutilise la même paire de cuivre en exploitant des fréquences que la voix n'utilise pas. La fibre optique, techniquement mature depuis les travaux de Corning en 1970 et le premier câble transatlantique TAT-8 en 1988, ne remplace la boucle locale cuivre que progressivement, parce que le cuivre est déjà posé et que la fibre exige de creuser de nouvelles tranchées. Pendant toute la période de transition, les opérateurs maintiennent la double infrastructure : les anciens commutateurs de circuits \emph{et} les nouveaux routeurs IP. La transition numérique est un double coût d'infrastructure, pas une substitution.
 
\dimmark{D6}
En 1989, Tim Berners-Lee propose au CERN un système d'hypertexte pour faciliter l'accès à l'information interne de l'organisation. Les protocoles qu'il développe --- HTTP, HTML, URL --- sont rendus publics en 1993. L'apparition du navigateur graphique Mosaic la même année, puis de Netscape en 1994, ouvre le Web au grand public. En 1998, Google introduit un moteur de recherche fondé sur l'analyse des liens hypertextes (le PageRank). Le nombre d'utilisateurs d'Internet dans le monde passe de quelques millions en 1993 à 280 millions en 1999 \parencite{abbate_inventing_1999}.
 
Le Web ajoute une dimension nouvelle au traitement de l'information : le stockage distribué permanent. Chaque page web est un fichier hébergé sur un serveur allumé vingt-quatre heures sur vingt-quatre, accessible à tout moment par n'importe quel utilisateur connecté. C'est la première infrastructure de STORE permanent à l'échelle mondiale. Un livre peut être fermé sans coût ; un serveur web ne peut pas être éteint sans que les pages qu'il héberge disparaissent. La mise en réseau du stockage transforme un acte ponctuel (stocker un fichier sur un disque local) en une obligation continue (alimenter le serveur qui rend le fichier accessible).
 
\dimmark{0-TER}
C'est dans cette période que l'acronyme \emph{ICT} (\emph{Information and Communication Technology}) se stabilise. L'OCDE et l'UNESCO l'adoptent à la fin des années 1990 pour désigner la convergence des télécommunications, de l'informatique et de l'audiovisuel. L'acronyme consacre l'unification de trajectoires techniques jusque-là distinctes. Il entérine aussi, dans son silence sur l'énergie, l'occultation que cette thèse cherche à lever : \emph{ICT} ne contient ni \emph{E} pour \emph{energy} ni \emph{M} pour \emph{materials}. La dimension 0-TER est absente du nom même du secteur.
 
La matérialité de l'infrastructure de réseau reproduit les structures observées au XIXe~siècle pour le télégraphe. Les câbles sous-marins en fibre optique suivent largement les routes des câbles télégraphiques posés un siècle plus tôt. La silice ultrapure nécessaire à la fabrication de la fibre, les navires câbliers spécialisés, les stations de répétition sous-marines sont invisibles pour l'utilisateur, tout comme l'étaient la gutta-percha malaise et les poteaux télégraphiques du réseau Chappe. Le parallèle matériel avec §\ref{subsec:materialite} est structurel, pas métaphorique : même fonction (TRANSMIT longue distance), même occultation du substrat physique, même concentration géographique de la production.

\dimmark{0-GOV/D10}{+}%
La commercialisation d'Internet entre 1993 et 2000 n'efface pas l'existence de bifurcations institutionnellement portées et politiquement écartées, en symétrie avec les architectures alternatives bloquées de la section précédente (§\ref{subsec:alternatives_bloquees}). Trois moments le rappellent. En 1971, AT\&T se voit offrir par l'ARPA la possibilité de racheter l'infrastructure ARPANET~; après examen, l'entreprise refuse, jugeant que la commutation par paquets ne sera jamais commercialement viable~\parencite{rose_att_2012, abbate_inventing_1999}. La décision laisse l'ARPANET sous tutelle militaire puis universitaire jusqu'à la privatisation de NSFNET en 1995, et ouvre l'espace dans lequel les fournisseurs d'accès commerciaux des années 1990 viendront se loger~; \textcite{russell_open_2014} a documenté plus largement le scepticisme des opérateurs télécommunications historiques à l'égard de la commutation par paquets pendant la majeure partie des années 1970, scepticisme qui a paradoxalement permis l'émergence d'Internet en dehors de leur tutelle. En 1994, le sénateur démocrate Daniel Inouye propose un projet de loi qui aurait obligé les opérateurs de télécommunications à réserver vingt pour cent de leurs capacités à des usages publics, éducatifs et non commerciaux~\parencite{tarnoff_internet_2022}~; le projet ne passe pas, faute d'un mouvement social organisé pour le porter. La privatisation intégrale de NSFNET suit en 1995, et le \emph{Framework for Global Electronic Commerce} promulgué par l'administration Clinton en juillet 1997 consacre cette orientation en posant l'autorégulation industrielle comme doctrine internationale du numérique~\parencite{clinton_framework_1997}. Ces trois moments ne sont pas des accidents~: ils marquent les points où une autre trajectoire d'Internet, plus régulée et plus partagée, aurait été institutionnellement possible mais a été écartée, et leur cumul fixe pour les vingt années suivantes un cadre où l'autorégulation marchande paraît la seule architecture envisageable.


% -----------------------------------------------------------------------------
\subsection{Le capital intangible et la mesure de l'invisible}\label{subsec:logiciel}
% -----------------------------------------------------------------------------
 
\dimmark{D3}
La séparation du hardware et du software imposée à IBM en 1969\footnote{Sous la pression de l'enquête antitrust ouverte par le Department of Justice (\emph{United States v.~IBM}, lancée en janvier 1969), IBM annonce dès le 23~juin 1969 sa nouvelle politique de \emph{unbundling}~: facturer séparément le hardware, le software et les services, ce qui ouvre la voie à une industrie indépendante du logiciel. L'enquête antitrust elle-même s'étend jusqu'en janvier 1982, date à laquelle elle est abandonnée sans condamnation par l'administration Reagan.} a créé les conditions d'une industrie du logiciel indépendante. En trente ans, cette industrie passe d'une activité artisanale de programmation sur mesure à un marché mondial dominé par quelques acteurs (Microsoft pour les systèmes d'exploitation, Oracle pour les bases de données, SAP pour la gestion intégrée). Le logiciel cesse d'être un poste de R\&D interne pour devenir un bien marchand, puis une catégorie de capital dans les comptes nationaux.
 
\textcite{corrado_measuring_2005} et \textcite{corrado_intangible_2009} fournissent la mesure de cette transformation. Leur estimation montre que l'investissement total des entreprises américaines dans le capital intangible --- logiciel, R\&D, design, formation, capital organisationnel, image de marque --- atteignait environ mille milliards de dollars en 1999, un montant approximativement égal à l'investissement en capital tangible la même année. Le basculement se situe au milieu des années 1990. En 2003, environ huit cent milliards de dollars d'investissement intangible restaient exclus des comptes nationaux publiés, ce qui conduit à l'exclusion de plus de trois mille milliards de dollars de stock de capital intangible.

\dimmark{D8/D6}{+}%
La croissance du capital intangible n'est pas uniformément distribuée à travers les secteurs. Un acteur peu visible dans la chronologie populaire de l'internet en a pourtant été l'un des premiers déployeurs massifs~: le secteur financier. \textcite{schiller_digital_2014} documente avec une précision quantitative rare que l'\emph{American Bankers Association} reconnaît en 1970 les réseaux de données comme \emph{operational basis for finance}, quinze ans avant la diffusion grand public de la téléinformatique. Le budget de traitement de données de Morgan Stanley est multiplié par cinq entre 1979 et 1985~; les premières formes d'\emph{algorithmic trading} apparaissent au début des années 1980 et se généralisent dans les deux décennies suivantes~\parencite{mackenzie_trading_2021}. La financiarisation des économies occidentales au sens que lui donne \textcite{epstein_financialization_2005} est techniquement portée par le déploiement précoce des ICT dans la banque d'investissement, le trading et la gestion d'actifs, bien avant que la commercialisation d'internet ne rende ce déploiement visible. Cette antériorité est analytiquement importante~: une part du capital intangible que mesurera \textcite{corrado_intangible_2009} dans les années 1990 a déjà été investie par le secteur financier dans la décennie précédente, et l'angle économique sous lequel le secteur ICT sera lu à partir de 2000 est en grande partie un héritage de l'angle financier sous lequel il a été déployé entre 1970 et 1990.

\dimmark{D1}
L'inclusion du capital intangible dans le cadre de la comptabilité de la croissance modifie les résultats observés. \textcite{corrado_intangible_2009} montrent que le taux de croissance de la production par travailleur augmente plus rapidement lorsque les actifs intangibles sont comptabilisés, et que l'approfondissement du capital (\emph{capital deepening}) devient la source dominante sans ambiguïté de la croissance de la productivité du travail. Le capital intangible n'est donc pas un détail comptable ; c'est la catégorie qui détermine l'interprétation des sources de la croissance économique dans la période 1995--2007.
 
\dimmark{0-TER}
Le paradoxe que ces mesures ne posent pas est le suivant : le capital intangible est traité dans les comptes nationaux comme s'il n'avait pas de substrat matériel. Le logiciel ne s'use pas physiquement (pas d'amortissement par usure) ; il n'a pas de localisation contrainte (il peut être copié) ; il ne consomme pas d'énergie en tant que tel (une ligne de code sur un disque n'a pas de coût énergétique propre). Mais le logiciel n'existe que s'il est exécuté par un processeur alimenté en électricité, et les données qu'il traite n'existent que sur des serveurs alimentés en permanence. Le capital intangible a un coût énergétique indirect que la catégorie comptable rend invisible. La séparation entre « intangible » et « tangible » dans la comptabilité nationale reproduit, au niveau des comptes, la séparation entre économie de l'information et économie de l'énergie que cette thèse identifie au niveau disciplinaire.

\dimmark{D3/0-TER}{*}%
Ce découplage comptable a une conséquence pratique~: les entreprises qui investissent massivement dans le capital intangible (logiciels, données, algorithmes) externalisent la consommation énergétique correspondante vers les fournisseurs de cloud et les opérateurs de datacenters, dont les coûts apparaissent comme des charges d'exploitation et non comme l'amortissement d'un actif. La structure comptable rend donc structurellement difficile l'imputation du coût énergétique de l'information à ses bénéficiaires économiques. Le chapitre~2 reprendra ce mécanisme sous la dimension D3~; le chapitre~3 tentera d'en estimer l'ampleur.

\questionch{La mesure du capital intangible de \textcite{corrado_intangible_2009} peut-elle être complétée par une estimation du coût énergétique indirect de ce capital ? Le chapitre~3 posera cette question en croisant les séries de capital intangible avec les estimations de consommation d'énergie du secteur ICT de \textcite{malmodin_energy_2018}.}
 
 
% -----------------------------------------------------------------------------
\subsection{La bulle spéculative et le cycle Perez}\label{subsec:dotcom}
% -----------------------------------------------------------------------------
 
\dimmark{D8}
Le NASDAQ passe d'environ mille points en 1995 à 5\,048 le 10 mars 2000, puis retombe à 1\,114 en octobre 2002. La capitalisation détruite entre le sommet et le creux dépasse cinq mille milliards de dollars. L'épisode n'est pas un accident ; c'est une régularité structurelle que \textcite{perez_technological_2002} formalise dans un modèle qui distingue cinq phases dans la relation entre capital financier et révolution technologique : irruption, ferveur spéculative (\emph{frenzy}), point de retournement (\emph{turning point}), synergie et maturité.
 
Le même schéma se rejoue pour chaque GPT. La canal mania des années 1790 en Grande-Bretagne, la railway mania des années 1840, la corporate mania électrique des années 1880--1890 et le radio boom des années 1920 précèdent chacun un krach suivi d'une phase de déploiement systématique de l'infrastructure. Le krach ne détruit pas l'infrastructure ; il la transfère à bas coût des spéculateurs vers les utilisateurs. Les chemins de fer construits pendant la bulle des années 1840 sont les mêmes qui transportent les marchandises dans les années 1850 ; les câbles de fibre optique posés pendant la bulle dotcom sont les mêmes qui portent le trafic Internet dans les années 2010.
 
L'épisode Iridium condense le même cycle à l'échelle d'un projet. Motorola investit cinq milliards de dollars dans une constellation de soixante-six satellites de communication en orbite basse, lancée commercialement en novembre 1998. Neuf mois plus tard, la société dépose le bilan, faute de clients (le prix du terminal est de 3\,000 dollars, celui de la minute de communication de 3 à 8 dollars, à un moment où le téléphone cellulaire terrestre devient bon marché et couvre l'essentiel du marché solvable). La constellation est rachetée en décembre 2000 pour 25 millions de dollars, soit 0,5~\% de l'investissement initial. L'infrastructure survit à la faillite de son promoteur et fonctionne encore aujourd'hui pour des usages de niche (maritime, aviation, zones non couvertes). Le parallèle avec les faillites des premiers câbles transatlantiques de 1858 est exact.

\dimmark{D5/D6}{+}%
La bulle laisse cependant une tension irrésolue que le seul cycle Perez ne couvre pas. La phase \emph{frenzy} a vu se déployer une infrastructure et une base d'utilisateurs sans qu'un modèle commercial stable n'émerge~: les abonnements payants sont massivement rejetés par les utilisateurs hérités de la culture du partage gratuit (§\ref{subsec:pc}), et la publicité bannière des premiers sites reste peu ciblée et peu rentable. Une formulation conceptuelle anticipe la résolution sans la produire~: \textcite{goldhaber_attention_1997} publie en avril 1997 dans \emph{First Monday} un article intitulé «~The Attention Economy and the Net~» dans lequel il propose que la ressource désormais rare et donc objet de la concurrence économique n'est pas l'information~--- abondante et reproductible à coût marginal nul~--- mais l'attention de l'utilisateur. La formulation reste prospective en 1997~: aucun modèle commercial ne l'opérationnalise au moment où elle est écrite. Mais elle nomme avec précision ce que les acteurs de l'industrie publicitaire et des moteurs de recherche cherchent sans le formuler~--- un dispositif qui capture et valorise l'attention plutôt que l'information. La résolution arrivera trois ans plus tard avec l'invention par Google d'un dispositif qui ne se contente pas de capturer l'attention mais l'inscrit dans un marché de prédictions comportementales~: mécanisme dont la section~\ref{sec:contemporain} examinera la portée et les conséquences matérielles.

\emergence{Le cycle Perez montre que le financement de l'infrastructure informationnelle obéit aux mêmes dynamiques que celui des infrastructures physiques antérieures. Le capital financier afflue vers une technologie dont la valeur perçue dépasse la valeur d'usage immédiate, crée une bulle, s'effondre, et laisse derrière lui une infrastructure que la phase de déploiement exploite. Ce résultat confirme que l'ICT n'est pas immatérielle : elle exige des investissements en capital fixe (câbles, serveurs, antennes, satellites) dont la logique financière est celle de toute infrastructure lourde.}
 
 
% -----------------------------------------------------------------------------
\subsection{Le téléphone intelligent et la convergence des fonctions}\label{subsec:smartphone}
% -----------------------------------------------------------------------------
 
\dimmark{0-TER}
Les données de \textcite{hilbert_worlds_2011} montrent qu'en 2007, la part numérique atteint 99,9~\% pour la télécommunication, 94~\% pour le stockage, mais seulement 25~\% pour le broadcasting. La convergence est donc achevée pour deux fonctions sur trois, et en cours pour la dernière. C'est dans cette configuration que l'iPhone d'Apple, lancé en juin 2007, rend la convergence visible dans un objet unique.
 
Un seul appareil réunit un processeur ARM (COMPUTE), une mémoire flash (STORE), un réseau cellulaire et Wi-Fi (TRANSMIT bidirectionnel), un écran tactile (interface), un récepteur GPS, un accéléromètre et une caméra (capteurs). Le routage IP est du calcul ; le cache du navigateur est du stockage ; le protocole HTTP est du contrôle. Les trois fonctions sont techniquement inséparables dans l'architecture de l'appareil. Ce que la tripartition TRANSMIT/STORE/COMPUTE distingue analytiquement, le smartphone fusionne physiquement.
 
\dimmark{D2}
Le smartphone n'invente pas la convergence ; il la rend portable et permanente. L'appareil est allumé en continu, connecté en permanence, et génère un flux constant de données (géolocalisation, consultation de pages, échanges de messages) qui transite par une infrastructure de réseau elle aussi permanente. Le résultat structurel est un déplacement du COMPUTE et du STORE de l'appareil vers le serveur distant : ce que le PC des années 1980 avait distribué chez l'utilisateur, le smartphone le recentralise dans le datacenter. L'architecture \emph{client léger / serveur lourd} qui en résulte est le point de départ de la section suivante (§\ref{sec:datacenter}).
 
Les données de la Banque mondiale, issues de l'UIT, confirment que le basculement se situe autour de 2005 : les abonnements mobiles dépassent les lignes fixes cette année-là (33,9 pour cent habitants contre 19,1), et le nombre de lignes fixes commence à décliner, passant de 19,1 en 2005 à 10,6 en 2023. Le mobile atteint 109,8 abonnements pour cent habitants en 2023, soit plus d'un abonnement par personne sur Terre. Le téléphone fixe, technologie dominante pendant un siècle (§\ref{subsec:telephone}), devient en vingt ans un résidu.
 
\emergence{La convergence T/S/C sur un substrat numérique unique est le résultat structurel de la période 1975--2007. Elle produit trois conséquences pour la thèse. Premièrement, la consommation d'énergie de l'information cesse d'être ponctuelle pour devenir permanente : le réseau est allumé en continu, les serveurs fonctionnent en permanence, le smartphone est toujours connecté. Deuxièmement, les demandes énergétiques des trois fonctions se composent au lieu de rester séparées : le TRANSMIT numérique contient du COMPUTE et du STORE, le COMPUTE contient du STORE et du TRANSMIT. Troisièmement, la mesure économique du secteur (le capital intangible de \textcite{corrado_intangible_2009}, le PIB du secteur ICT de l'OCDE, les indices de prix hédoniques du BEA) reste séparée de la mesure énergétique (\textcite{koomey_implications_2011}, \textcite{malmodin_energy_2018}, IEA). Personne ne croise les deux. C'est dans cet écart que s'ouvre la section suivante.}
 
\questionch{Comment formaliser le couplage entre capacités informationnelles (mesurées par \textcite{hilbert_worlds_2011} en bits et MIPS) et consommation d'énergie (mesurée par \textcite{koomey_implications_2011} en kWh par calcul) ? Le croisement des deux séries donnerait une estimation de la consommation d'énergie totale du traitement de l'information mondiale. Ce croisement n'a jamais été publié. Le chapitre~3 de cette thèse le tentera.}


\begin{table}[htbp]
\caption{Dimensions économiques de la convergence (§\ref{sec:convergence}).}\label{tab:convergence-dimensions}
\centering\small
\renewcommand{\arraystretch}{1.15}
\begin{tabularx}{\textwidth}{@{}cXc@{}}
\toprule
\textbf{Dim.} & \textbf{Ce que §\ref{sec:convergence} établit} & \textbf{Niveau} \\
\midrule
D1 & L'information numérique est non rivale~; le circuit qui la transporte est rival. Le paradoxe de la productivité (Solow 1987, David 1990) reflète un décalage de substrat, pas seulement d'organisation. & + \\
D2 & Plateformes \emph{winner-take-all} (Microsoft, Google)~; architecture ouverte IBM-PC $\rightarrow$ clonage $\rightarrow$ concentration sur l'OS et les applications. & + \\
D3 & Capital intangible $\geq$ tangible dès 1995 (\textcite{corrado_intangible_2009})~; \emph{unbundling} IBM (1969) comme origine institutionnelle~; découplage comptable entre valeur du logiciel et coût énergétique de son usage. & ★ \\
D5 & Paradoxe de la productivité (Solow 1987, David 1990)~; le gain de productivité n'apparaît qu'après recomposition organisationnelle complète. & ★ \\
D6 & Le Web (1993) crée un signal de prix instantané à l'échelle mondiale~; le moteur de recherche (Google, 1998) devient l'interface d'accès à l'information. & + \\
D7 & Standards ouverts (TCP/IP 1983, HTTP 1993) comme condition de la convergence~; modularité au sens de \textcite{langlois_modularity_2003}. & ★ \\
D8 & Cycle Perez~: dotcom = railway mania~; l'infrastructure (câbles, serveurs) survit au krach et sert la phase de déploiement. Secteur financier comme premier déployeur massif (1970--1990). & ★ \\
D10 & Open source (GNU 1984, Linux 1991) vs licensing massif (Microsoft)~: deux trajectoires coexistantes~; culture \emph{Whole Earth Catalog} comme trace structurelle. Bifurcations écartées (refus AT\&T 1971, échec Inouye 1994). & + \\
\midrule
0-TER & STORE devient permanent (serveur web \emph{always-on})~; premier régime de consommation continue~; la numérisation du réseau (1983--2000) est un double coût d'infrastructure, pas une substitution. & ★ \\
0-GOV & Privatisation NSFNET (1995)~; \emph{Framework for Global Electronic Commerce} (Clinton 1997) consacre l'autorégulation~; trajectoires alternatives (Inouye 1994) écartées. & + \\
\bottomrule
\end{tabularx}
\renewcommand{\arraystretch}{1.0}
\end{table}
